\documentclass{article}
\usepackage{amsmath, amssymb, bm}

\begin{document}

\section*{Problem 002: Partially Linear Single Index Model}

Let the observed data be i.i.d. copies of
\[
O=(Y,X,Z),
\]
where \(Y\in\mathbb R\), \(X\in\mathbb R^p\), and \(Z\in\mathbb R^q\).

Consider the partially linear single index model
\[
Y = X^\top \beta + g(Z^\top \alpha) + \varepsilon,
\qquad
E(\varepsilon\mid X,Z)=0,
\]
where \(\beta\in\mathbb R^p\), \(\alpha\in\mathbb R^q\), and
\(g:\mathbb R\to\mathbb R\) is an unknown smooth nuisance function.

Let
\[
\theta=(\beta^\top,\alpha^\top)^\top.
\]

\subsection*{Identifiability}

Because \(g(Z^\top\alpha)\) is invariant to scale changes in \(\alpha\), impose
the normalization
\[
\|\alpha\|_2=1,
\qquad
\alpha_1>0.
\]

Equivalently, one may use a local unconstrained parameterization
\[
\alpha=\alpha(\gamma),
\qquad
\gamma\in\mathbb R^{q-1},
\]
and define the finite-dimensional parameter as
\[
\vartheta=(\beta^\top,\gamma^\top)^\top.
\]

In the derivation, please state clearly whether the influence function is for
\(\theta=(\beta^\top,\alpha^\top)^\top\) under the constraint, or for the
unconstrained local parameter \(\vartheta=(\beta^\top,\gamma^\top)^\top\).

\subsection*{Density factorization}

The joint density may be written as
\[
f_{X,Z,Y}(x,z,y)
=
f_{X,Z}(x,z)
f_{\varepsilon\mid X,Z}
\left\{
y-x^\top\beta-g(z^\top\alpha),x,z
\right\},
\]
with the conditional mean restriction
\[
E(\varepsilon\mid X,Z)=0.
\]

\subsection*{Question}

Derive the influence function or efficient influence function for the
finite-dimensional parameter of interest in this semiparametric model.

Please:

\begin{enumerate}
\item normalize the problem and state the target parameter precisely;
\item state the nuisance functions and model restrictions;
\item derive the relevant score or moment equations;
\item characterize the nuisance tangent space;
\item derive a candidate influence function for \(\beta\) and \(\alpha\), or
      for \(\beta\) and the local parameter \(\gamma\);
\item state the projection needed for semiparametric efficiency;
\item clearly mark any unresolved steps if a full EIF derivation requires
      additional smoothness, completeness, or invertibility assumptions.
\end{enumerate}

\end{document}

